%! Systems Involving Quadratics — Unit 3, Lesson 3.6 — Homework
\documentclass[10pt]{article}
\usepackage{algebra2-article}
\usepackage{algebra2-boxes}
\newcommand{\TallMath}[1]{$\displaystyle #1\rule[-1.4em]{0pt}{3.2em}$}

\begin{document}

\pageheader{Unit 3, Lesson 3.6}{Homework}

\vspace{0.10in}

\begin{notesbox}{Practice}
Show your work. Every solution is an \textbf{ordered pair} $(x,y)$ --- find both coordinates and check them in
\emph{both} equations.

\begin{enumerate}[leftmargin=1.9em, itemsep=9pt]
  \item \textbf{Is the pair a solution?} Check each against $\ y=x^2+2\ $ and $\ y=2x+1$ (write yes/no).
    \[ (1,3):\ \blank{1.6cm} \qquad\qquad (0,2):\ \blank{1.6cm} \]

  \item \textbf{Solve by substitution} (linear--quadratic; real, rational). Write ordered pairs.
    \begin{enumerate}[label=(\alph*), leftmargin=1.9em, itemsep=4pt]
      \item $y=x^2-2x-3,\ \ y=2x-3$
        \begin{work}
          x^2-2x-3 &= 2x-3 \\
          x^2-4x &= 0 \\
          x(x-4) &= 0 \\
          x &= 0,\ 4 \\
          y &= -3,\ 5 \quad\text{so } (0,-3) \text{ and } (4,5)
        \end{work}
      \item $y=x^2-4x+3,\ \ y=x-1$
        \begin{work}
          x^2-4x+3 &= x-1 \\
          x^2-5x+4 &= 0 \\
          (x-1)(x-4) &= 0 \\
          x &= 1,\ 4 \\
          y &= 0,\ 3 \quad\text{so } (1,0) \text{ and } (4,3)
        \end{work}
    \end{enumerate}

  \item \textbf{Won't factor} --- use the quadratic formula. Give the $x$-coordinates in simplest radical form.
    \begin{work}
      x^2+1 &= 3x \\
      x^2-3x+1 &= 0 \\
      x &= \frac{3\pm\sqrt{(-3)^2-4(1)(1)}}{2(1)} \\
        &= \frac{3\pm\sqrt{5}}{2}
    \end{work}

  \item \textbf{Quadratic--quadratic.} Set the expressions equal.
    \begin{enumerate}[label=(\alph*), leftmargin=1.9em, itemsep=4pt]
      \item $y=x^2,\ \ y=-x^2+18$
        \begin{work}
          x^2 &= -x^2+18 \\
          2x^2 &= 18 \\
          x &= \pm3 \\
          y &= 9 \quad\text{so } (-3,9) \text{ and } (3,9)
        \end{work}
      \item $y=x^2+2x,\ \ y=x^2-4$ \quad\emph{(watch the $x^2$ terms!)}
        \begin{work}
          x^2+2x &= x^2-4 \\
          2x &= -4 \\
          x &= -2 \\
          y &= 0 \quad\text{so } (-2,0) \text{ only}
        \end{work}
    \end{enumerate}

  \item \textbf{How many solutions?} Each system collapses to the quadratic shown. Compute $b^2-4ac$ and state how
        many points the graphs share ($0$, $1$, or $2$).
    \[ x^2-4x+4=0:\ \blank{2.6cm} \qquad x^2-x-6=0:\ \blank{2.6cm} \qquad x^2+x+3=0:\ \blank{2.6cm} \]

  \item \textbf{Find the error.} Priya solves $\ y=x^2,\ \ y=x+6\ $ and correctly gets $x=3$ and $x=-2$. She writes
        ``the solutions are $x=3$ and $x=-2$.'' What did she forget, and what are the correct solutions?
        \writeline
\end{enumerate}
\end{notesbox}

\begin{extensionbox}
\textbf{1. Model (break-even).} A start-up's revenue is $y=-x^2+12x$ and its cost is $y=4x+12$ (hundreds of
dollars, $x$ in hundreds of units). Find both break-even points and the interval where it turns a profit.
\[ x=\blank{2.4cm} \quad\text{points: } \blank{3.2cm} \quad\text{profit when } \blank{2.4cm} \]

\textbf{2. Make it tangent.} For what value of $c$ is the line $y=4x+c$ tangent to $y=x^2$ (exactly one solution)?
\emph{Set equal, then force $b^2-4ac=0$.} \blank{2.0cm}
\end{extensionbox}

\begin{spiralbox}
\textbf{Coming next --- Lesson 3.7: Modeling with Quadratics.} You've now solved single quadratics \emph{and}
systems. Next you'll build the quadratic model yourself from a projectile, area, or optimization story --- and read
its \emph{vertex} as the maximum height, maximum area, or best price. The algebra is all from this unit; the new
skill is translating the situation into the model and interpreting the answer.
\end{spiralbox}

\end{document}
